% Tapis de Quillen, with Quillen's offprint and preprint (cote 162-5, typed part): archive cote 162-5, site page 37 (PDF page 38).
% typescript. Draft compared: none.
% AI-transcribed by Claude Opus 5.5 (claude-opus-5-5, adaptive thinking, effort high), Message Batches API, batch t162_full (msgbatch_01UAh3zyWPd454AeDkhnmQ33), 2026-11-13; not yet checked by eye.
\scanpage{162-5/37}
morphismes de flêches.

Notons une symmétrie évidente dans la c atégorie (Bicat), interprétée
en termes d\xout{\unsure{e}}'ensemble bi-semisimpliciaux, obtenue en échangeant les lignes
et les colonnes dans \xout{les} de tels ensembles bi... Un ensemble bisemisimplici-%
al qui ... peut donc de deux façons différentes être considéré comme défi-%
nissant une bicatégorie. De même le groupe symétrique $\mathrm{Sym}_n$ opère sur
la catégorie \xout{de} \unsure{(nupcat)|(n-upcat)} des catégories n-uples ..

Une 2-catégorie peut être définie comme constituée par un ensemble
\xout{discret} $C^{\hand{\unsure{0}}}_{\hand{0}}$ (l'ensemble des objets ), qu'on interprètera comme une catégo
rie discrète, plus la donnée pour tout x,y \hand{$\in$}$C^{\hand{0}}_{\hand{0}}$ d'une catégorie $\underline{\mathrm{Hom}}$(x,y),
plus \inserted[typed]{des foncteurs} \xout{des} composition entre celles\unsure{\hand{-}}ci ... Prenant la catégorie somme \xout{de} $C^{\hand{1}}_{\hand{1}}$
des $\underline{\mathrm{Hom}}$(x,y), On peut exprimer la donnée de ces foncteurs de composition,
et des sources et buts \inserted[typed]{x,y} dans les catégories $\underline{\mathrm{Hom}}$(x,y), comme définis par
des foncteurs ``source'' et ``but'' $C^{\hand{1}}_{\hand{1}} \hand{\to} C^{\hand{0}}_{\hand{0}}$, et un foncteur de composition
$C^{\hand{1}}_{\hand{1}} \mathrm{x}_{C^{\hand{0}}_{\hand{0}}} C^{\hand{1}}_{\hand{1}} \hand{\longrightarrow} C^{\hand{1}}_{\hand{1}}$ , les axiomes des 2-catégories correspondant exactement aux
axiomes des(Cat)-catégories. De cette façon, la catégorie des 2-catégories
est équivalente à la catégorie des (Cat)-catégories qui ont une ``catégorie
d'objets'' qui est discrète, ou encore à la catégorie des Bicatégories
qui, interprétées comme ensembles bi-semisimpliciaux, ont \xout{une} comme
première ligne un ensemble semisimplicial discret (i.e. un foncteur cons-%
tant sur \hand{$\Delta$} ). Il parait qu'on a une interprétation simple dans le même
esprit de la notion de n-catégorie en termes d'ensembles n-semi-simpliciaux
\margin{$\|$} \xout{xx} Demander à Quillen relations entre la notion de m-Catégorie dans
\margin{$\|$} la catégorie (n-Cat), et celle de ensemble (m+n)-semisimplicial.
\margin{$\|$} \xout{(en particulier entre cette dernière notion et celle de}

L'interprétation standard semisimpliciale permet d'interpréter

\mspage{\hand{37}}
